\documentclass[11pt]{article}
\usepackage{preamble}
\renewcommand{\answer}[1]{}

\begin{document}

\ladderheading{AP Statistics \textperiodcentered\ Unit 1 \textperiodcentered\ Topic 1.9 \textperiodcentered\ B: 2026-09-24 \textperiodcentered\ E: 2026-09-30}{Which Class Did Better?}

\focusbanner{TODAY'S PROBLEM}

Two AP Statistics classes took the same 40-point quiz. The side-by-side boxplots show the scores. Two students read them differently. \textbf{Sam:} ``Period B did better --- its highest score is 40 and Period E's is only 38.'' \textbf{Tia:} ``Period E did better --- its median is higher and its scores are more consistent.''\par
\begin{center}
\begin{tikzpicture}
\begin{axis}[
  width=0.85\linewidth,
  height=2.00in,
  ytick={2,1}, yticklabels={{Period B},{Period E}}, ymin=0.4, ymax=2.6,
  axis x line=bottom, axis y line*=left, y axis line style={draw=none}, boxplot/box extend=0.5, xmajorgrids, enlarge x limits=0.06,
  xlabel={Quiz score (out of 40)}, tick label style={font=\small}, label style={font=\small},
  ytick style={draw=none}, yticklabel style={font=\small\bfseries}
]
\addplot[boxplot prepared={lower whisker=12, lower quartile=22, median=27, upper quartile=33, upper whisker=40, draw position=2}, fill=calloutblue, draw=dayblue] coordinates {};
\addplot[only marks, mark=*, color=warmred] coordinates {(5,2)};
\addplot[boxplot prepared={lower whisker=18, lower quartile=26, median=31, upper quartile=34, upper whisker=38, draw position=1}, fill=calloutblue, draw=dayblue] coordinates {};
\end{axis}
\end{tikzpicture}
\end{center}
\begin{firsttakebox}{FIRST TAKE --- before the video (5 min)}
Which class did better on the quiz? One sentence, and name ONE feature of the plot you used.
\workspace{0.9in}
\end{firsttakebox}

\begin{callout}{\IconBook\ COMPARING DISTRIBUTIONS (keep on the page)}{calloutgreen}
Use comparative words such as higher, lower, more variable, or similar. Compare center (median), variability ($\mathrm{IQR}=Q_3-Q_1$, range $=$ maximum minus minimum), apparent shape, and unusual features in context. Include plotted outliers when finding the overall range. A boxplot hides values within each quartile, so qualify shape claims with `appears.' A median describes center; a maximum describes the highest observed value.
\end{callout}

\newpage
\textbf{Finish after the video:}\par\smallskip

\dokbadge{1}~\textbf{(a)} Read the median for each class.\par
\workspace{0.6in}

\dokbadge{2}~\textbf{(b)} In two or three sentences, compare the quiz scores' center, spread, apparent shape, and unusual features. Use points and comparative words.\par
\workspace{1.5in}

\dokbadge{3}~\textbf{(c)} Whose claim is better supported, Sam's or Tia's? Cite two specific features of the boxplots that settle it. Then: if you could report only ONE statistic per class to a parent asking `how did the class do?', which statistic would you choose, and why that one rather than the maximum?\par
\workspace{2.0in}

\begin{sentenceframebox}\raggedright\textbf{Frame:} \blank[0.8in]'s claim is better supported because the \blank[1in] of Period E is \blank[0.8in] than Period B's, and Period E's scores are \blank[1in].\end{sentenceframebox}

\begin{sentenceframebox}\raggedright\textbf{Frame:} I would report the \blank[0.9in] because the maximum describes \blank[1.6in].\end{sentenceframebox}

\begin{summaryexitbox}{EXIT --- turn this in}
One thing the video changed about my first take: \hrulefill\par
\workspace{0.4in}
\end{summaryexitbox}

\bankitem{Optional \#1 --- not collected}{\dokbadge{2}~Two bus routes were timed for 20 mornings each. Route 1: mean 24.0 min, SD 6.5 min, median 22 min. Route 2: mean 24.5 min, SD 1.8 min, median 24 min. Write two sentences comparing the routes for a student who must not be late, using comparative language and the statistics given.}{}

\end{document}
